\section{Introduction}\label{sec:introduction}
The planar XY model describes unit vectors interacting through the cosine of their angle difference. Its transition is expected to differ sharply from an ordinary power-law critical point: on the massive side, the correlation length diverges exponentially in the inverse square root of the distance to criticality. We prove this asymptotic, with a single limiting constant, for the square-lattice model with its original cosine interaction.

For $R\in\N$, let $\Lambda_R=[-R,R]^2\cap\Z^2$ and let $E_R$ contain each unordered nearest-neighbor edge of $\Lambda_R$ once. At inverse temperature $b>0$, the free-boundary measure is
\begin{equation}\label{eq:spin-law}
 \frac{1}{Z_{R,b}}\exp\!\left(b\sum_{\{u,v\}\in E_R}
                    \cos(\theta_u-\theta_v)\right)
       \prod_{u\in\Lambda_R}\frac{d\theta_u}{2\pi},
 \qquad \theta_u\in\R/(2\pi\Z).
\end{equation}
There are no exterior edges. Write $e_1=(1,0)$ and define
\begin{equation}\label{eq:mass-definition}
 C_b(r)=\lim_{R\to\infty}\E^{\mathrm{free}}_{R,b}
       \cos(\theta_0-\theta_{re_1}),\qquad
 m(b)=\lim_{r\to\infty}-\frac1r\log C_b(r),\qquad
 b_c=\inf\{b>0:m(b)=0\}.
\end{equation}
The thermodynamic limit precedes the limit in $r$. We use the standard existence of these limits and the facts $0<b_c<\infty$, $m(b_c)=0$, and $0<m(b)<\infty$ for $0<b<b_c$; see the correlation inequalities and phase-transition results discussed below. The axis correlation length in the massive phase is $\xi(b)=m(b)^{-1}$.

\begin{theorem}\label{thm:main}
For the model \eqref{eq:spin-law}, there is a constant $A_{\mathrm{XY}}\in(0,\infty)$ such that
\begin{equation}\label{eq:main-limit}
 \lim_{b\uparrow b_c}\sqrt{b_c-b}\,\log\xi(b)=A_{\mathrm{XY}}.
\end{equation}
\end{theorem}
The temperature parameter in this limit is exactly the coefficient $b$ in \eqref{eq:spin-law}. The constant depends on this normalization and is not evaluated in closed form.

\subsection{Context and preceding results}
Berezinskii predicted a low-temperature phase with algebraically decaying correlations \cite{Berezinskii}. Kosterlitz and Thouless identified vortex binding as the mechanism of the transition \cite{KT}. Kosterlitz's renormalization analysis predicted the essential singularity of the correlation length \cite{Kosterlitz}. The quadratic flow underlying that prediction also underlies our argument. Its integration is explicit; the mathematical difficulty is to justify how the original temperature moves the microscopic model across the critical trajectory, and to recover the infinite-volume mass from the resulting scale.

Fr\"ohlich and Spencer established the low-temperature regime of slow decay for plane rotators and the associated vortex-binding mechanism \cite{FS}. The modern current and loop approach of van Engelenburg and Lis gives an elementary proof of the planar XY phase transition and a useful exponential-versus-polynomial correlation dichotomy \cite{vEL}. Their loop representation and finite-size criterion are used in our final passage from height fluctuations to spin correlations. The underlying correlation inequalities go back to Ginibre and to the Lieb--Rivasseau bounds for plane rotators \cite{Ginibre,Lieb,Rivasseau}.

Aizenman, Harel, Peled, and Shapiro developed a complementary height-level-set approach for the Villain and XY models, using in particular the lattice-Gaussian inequality of Regev and Stephens-Davidowitz \cite{AHPS,RSD}. Lammers proved an exact relation between spin and dual-height correlation lengths \cite[Theorem~1]{Lammers}. Durand and Lammers further express the XY mass through the right derivative at zero slope of the dual-height free energy \cite[Theorem~1.4]{DurandLammers}. The issue here is to control height fluctuations on scales diverging together with the temperature. Rigorous renormalization-group analyses include Falco's critical small-activity Coulomb gas and Bauerschmidt, Park, and Rodriguez's Gaussian scaling limits for the high-temperature discrete Gaussian model \cite{Falco,BPR1,BPR2}. These results concern different hypotheses from the ordinary cosine XY model at its fixed critical temperature.

Two companion works provide the critical inputs for the present proof. The first identifies the critical correlation exponent and the critical coefficient of the dual height field \cite{X}. The second establishes the logarithmic correction to the critical correlation and develops the analytic activity map used here \cite{F}. In the height normalization specified in Section~\ref{sec:finite}, the critical coefficient is $8\pi$. We use the companions' finite-graph comparisons, critical height limits, and analytic map estimates at their stated scope. The extension to a shrinking temperature neighborhood, the nonzero temperature response, and the conversion of a moving finite-scale profile into the near-critical mass are proved below.

\subsection{Proof strategy and new estimates}
The proof separates the determination of a scale from its interpretation as a correlation length. This separation allows the analytic activity map to be used only while its coordinates are small.

The first task is to make the temperature a controlled variable in that map. After averaging heights over cells of side $s$, we prove analytic initialization on a complex disk of radius proportional to $s^{-2}$ about $b_c$. A Poisson representation controls changes of physical weights on the finitely many comparison graphs. For complex parameters, we compare unnormalized physical factors to a nearby real law and retain the same quadratic regulator. Only local scalar normalizers are inverted. Interpolation between the real smallness estimate and a complex majorant then preserves the small activity bound.

At scale $L^j$, the two scalar coordinates $x_j$ and $y_j$ measure, respectively, a quadratic gradient perturbation and the first Fourier harmonic of the coarse height weight. After normalization, their leading recursion is
\begin{equation}\label{eq:intro-map}
 (x,y)\longmapsto(x-y^2,\ y-xy).
\end{equation}
The remaining coordinates contract. At criticality, $(x_j,y_j)=(1,1)/j+O(\log j/j^2)$. Differentiating in the original parameter $b$, we prove
\begin{equation}\label{eq:intro-response}
 \frac1j\,\partial_b(x_j,y_j)\big|_{b=b_c}
       \longrightarrow\mathcal P(1,-1/2),\qquad 0<\mathcal P<\infty.
\end{equation}
The positivity of $\mathcal P$ is essential. It follows from a finite-volume lower bound on the derivative of a height variance. Comparing common physical observables allows derivatives from different observation scales to be joined without identifying their different remainder spaces.

Put $\delta=\sqrt{b_c-b}$ and let $w=\delta j$ be the scaled logarithmic length. A comparison of finite differences on overlapping histories gives the initial condition for the rescaled coordinates $X=x/\delta$, $Y=y/\delta$ and their limiting flow
\begin{equation}\label{eq:intro-flow}
 X'=-Y^2,\qquad Y'=-XY,\qquad
 X(w)=k\cot(kw),\quad Y(w)=k\csc(kw),\quad k=\sqrt{3\mathcal P}.
\end{equation}
The flow reaches its outgoing pole at $T=\pi/k$. We prove that the discrete trajectory reaches any sufficiently small fixed outgoing level at an index $j_e$ satisfying $\delta j_e\to T$. Thus the relevant spatial scale is $L^{T/\delta}$, where $L$ is the fixed block ratio of the map.

The second task is to identify this scale with $\xi(b)$. Let $a_m(b)$ be the variance coefficient of an affine height test on a free square of side $m$. Along a sequence $b\uparrow b_c$, compactness and cutting comparisons produce a nonincreasing profile $\mathcal A(w)$ at scales $w=\delta\log_L m$. We establish Gaussian height limits in each continuity window, including the pinned geometries needed for spin duality. A relative growth estimate for a positive phase coefficient excludes intermediate profile values. The torus diagnostics before and at the exit then force
\begin{equation}\label{eq:intro-profile}
 \mathcal A(w)=8\pi\quad(0<w<T),\qquad
 \mathcal A(w)=0\quad(w>T).
\end{equation}
Figure~\ref{fig:flow-profile} relates the limiting flow to this height profile.

Below $T$, positive pinned-height covariance requires loops of macroscopic diameter and hence a correlation length of at least the corresponding exponential scale. Above $T$, the vanishing height coefficient makes fixed circulation tilts arbitrarily inexpensive. A stack of annuli gives arbitrarily strong polynomial decay at those scales; the Lieb--Rivasseau finite-size criterion then gives an upper bound on $\xi(b)$. Both bounds have the same limiting logarithmic scale. In particular,
\begin{equation}\label{eq:amplitude}
 A_{\mathrm{XY}}=(\log L)\frac{\pi}{\sqrt{3\mathcal P}}.
\end{equation}
Although the intermediate coordinates depend on the chosen block construction, the left side is determined by \eqref{eq:spin-law} and \eqref{eq:mass-definition}.

Section~\ref{sec:finite} develops the finite-graph estimates. Section~\ref{sec:entry} states the analytic map inputs and proves the temperature extension. Section~\ref{sec:trajectory} establishes \eqref{eq:intro-response} and the passage time. Section~\ref{sec:profile} identifies the moving-scale height profile, and Section~\ref{sec:mass} completes the transfer to the mass.
\begin{figure}[tb]
\centering
\begin{tikzpicture}[x=1cm,y=1cm,>=Latex,font=\small]
\begin{scope}[xshift=0cm,yshift=1.9cm,x=1.42cm,y=.39cm]
 \draw[->] (0,-4.2)--(0,4.7) node[above] {$X/k,\ Y/k$};
 \draw[->] (0,0)--(3.42,0) node[right] {$kw$};
 \draw[densely dashed,gray] (3.14159,-4.1)--(3.14159,4.3);
 \node[below right] at (3.14159,0) {$\pi$};
 \node[below] at (1.5708,0) {$\pi/2$};
 \draw[blue!65!black,very thick,domain=.26:2.88,samples=100]
  plot (\x,{cos(\x*180/pi)/sin(\x*180/pi)});
 \draw[orange!75!black,very thick,domain=.26:2.88,samples=100]
  plot (\x,{1/sin(\x*180/pi)});
 \node[blue!65!black,fill=white,inner sep=1pt] at (2.2,-2.8) {$\cot(kw)$};
 \node[orange!75!black,fill=white,inner sep=1pt] at (1.75,2.5) {$\csc(kw)$};
 \draw[->,blue!65!black,thick] (1.17,.43)--(1.34,.235);
 \node[align=center] at (1.57,-5.8) {Limiting scalar trajectory};
\end{scope}
\begin{scope}[xshift=7.2cm,yshift=.3cm,x=1.45cm,y=2.8cm]
 \draw[->] (0,0)--(3.4,0) node[right] {$w$};
 \draw[->] (0,0)--(0,1.35) node[above] {$\mathcal A(w)$};
 \draw[densely dashed,gray] (1.85,0)--(1.85,1.15);
 \node[below] at (1.85,0) {$T=\pi/k$};
 \draw[blue!65!black,very thick] (.06,1)--(1.85,1);
 \draw[blue!65!black,very thick] (1.85,.015)--(3.1,.015);
 \draw[blue!65!black,fill=white,thick] (1.85,1) circle[x radius=.045,y radius=.023];
 \draw[blue!65!black,fill=white,thick] (1.85,.015) circle[x radius=.045,y radius=.023];
 \node[left] at (0,1) {$8\pi$};
 \node[above,align=center,font=\footnotesize] at (.87,.09) {Positive pinned\\covariance};
 \node[above,align=center,font=\footnotesize] at (2.66,.35) {Vanishing\\height variance};
 \node[align=center] at (1.55,-.23) {Limiting height coefficient};
\end{scope}
\end{tikzpicture}
\caption{The two limits that determine the correlation length. Left: the explicit rescaled trajectory, with $k=\sqrt{3\mathcal P}$, has its outgoing pole at $w=T$. Right: the finite-scale free-height coefficient has a step at the same value. No value at the jump is asserted. Positive covariance below $T$ and circulation estimates above $T$ give the two bounds on $\delta\log_L\xi(b)$.}
\label{fig:flow-profile}
\end{figure}

